\renewcommand{\Anno}{2021}

\renewcommand{\NumeroProblema}{3}

\renewcommand{\Testo}{
Un obelisco di altezza $h>0$ è posizionato in un punto $P$, perpendicolarmente rispetto al terreno, che supponiamo piano ed infinito. A distanza $d>0$ da $P$ si trova una parete piana, anche questa perpendicolare al terreno, che supponiamo infinita sia in altezza sia in lunghezza. Sulla parete è posto un riflettore, che supponiamo puntiforme, la cui posizione può essere controllata dai tecnici. Nella notte, il riflettore illumina l'obelisco, che di conseguenza proietta la sua ombra sul terreno.

\begin{enumerate}
\renewcommand{\labelenumi}{(\alph{enumi})}

\item  Supponiamo che il riflettore si trovi in un punto della parete ad altezza~$4h$ rispetto al terreno.

Determinare quali posizioni può assumere sul terreno l'ombra della punta dell'obelisco.

\item  Posizioniamo ora il riflettore in modo che l'ombra dell'obelisco abbia una lunghezza uguale all'altezza dell'obelisco stesso, e sia perpendicolare rispetto alla parete.

Determinare come i tecnici devono muovere il riflettore sulla parete in modo che l'ombra della punta dell'obelisco si muova sul terreno a velocità costante verso la parete, restando su una retta che forma un angolo di~$45^{\circ}$ con la parete stessa.

\item  Supponiamo che il riflettore sia posizionato inizialmente come nella domanda precedente. Supponiamo che i tecnici vogliano fare in modo che l'ombra della punta dell'obelisco si muova come nella domanda precedente, ma lasciando fisso il riflettore e spostando la base dell'obelisco (lasciandolo ovviamente sempre perpendicolare al terreno).

Determinare come bisogna spostare la base dell'obelisco.

\end{enumerate}
}

\renewcommand{\Risposta}{
% SSSUP-BOZZA-AI: Risposta; autore=ChatGPT; data=2026-11-08; stato=DA-VERIFICARE
% SSSUP-PROVENIENZA: AI: elaborazione nuova dal testo
Poniamo la base iniziale in $(0,0)$, la parete nel piano $x=-d$ e usiamo $z$ per l'altezza.
\begin{enumerate}
\item[(a)] L'ombra della punta percorre tutta la retta del terreno $x=d/3$.
\item[(b)] Se la velocità desiderata è $v>0$, poniamo $s=vt/\sqrt2$. Per ottenere l'ombra $Q(t)=(h-s,s,0)$, il riflettore deve stare in
\[
L(t)=\left(-d,-\frac{ds}{h-s},h+\frac{hd}{h-s}\right),\qquad0\leq s<h.
\]
L'altra direzione a $45^\circ$ si ottiene cambiando segno alla coordinata $y$.
\item[(c)] Con il riflettore fisso in $(-d,0,h+d)$ basta muovere la base secondo
\[
P(t)=\left(-\frac{ds}{h+d},\frac{ds}{h+d},0\right).
\]
La base procede a velocità costante $vd/(h+d)$ nella stessa direzione dell'ombra.
\end{enumerate}
}

\renewcommand{\Aiutino}{
% SSSUP-BOZZA-AI: Aiutino; autore=ChatGPT; data=2026-11-08; stato=DA-VERIFICARE
% SSSUP-PROVENIENZA: AI: elaborazione nuova dal testo
Interseca con il piano del terreno la retta che passa per riflettore e punta dell'obelisco. Con obelisco fermo, l'ombra è una proiezione centrale; con riflettore fermo e altezza dell'obelisco costante, la posizione dell'ombra dipende invece in modo affine dalla posizione della base.
}

\renewcommand{\SoluzioneUfficiale}{}

\renewcommand{\SoluzioniAlternative}{}

\renewcommand{\Osservazioni}{}

\renewcommand{\Commenti}{}

% Aggiunta redazionale: la soluzione preesistente non viene modificata.
\providecommand{\SoluzioneRedazionale}{}
\renewcommand{\SoluzioneRedazionale}{
% SSSUP-BOZZA-AI: SoluzioneRedazionale; autore=ChatGPT; data=2026-11-08; stato=DA-VERIFICARE
% SSSUP-PROVENIENZA: AI: elaborazione nuova dal testo
Scegliamo coordinate con la base iniziale dell'obelisco in $P=(0,0,0)$, la punta in $U=(0,0,h)$ e la parete nel piano $x=-d$. Un riflettore $L=(-d,u,z)$ proietta un'ombra finita della punta se $z>h$. La semiretta da $L$ per $U$ incontra il terreno in
\[
Q=\left(\frac{hd}{z-h},-\frac{hu}{z-h},0\right).
\]

\paragraph{Punto (a).}
Per $z=4h$ si ottiene $Q=(d/3,-u/3,0)$. Al variare di $u$ si percorre tutta la retta $x=d/3$, parallela alla parete.

\paragraph{Punto (b).}
La configurazione iniziale richiesta ha ombra $Q(0)=(h,0,0)$; dalla formula precedente si ricava $L(0)=(-d,0,h+d)$. Una delle due traiettorie a $45^\circ$ verso la parete, percorsa con velocità $v$, è
\[
Q(t)=(h-s,s,0),\qquad s=vt/\sqrt2.
\]
Risolvendo rispetto a $u,z$ si trovano
\[
u=-\frac{ds}{h-s},\qquad z=h+\frac{hd}{h-s}.
\]
Il riflettore segue dunque sulla parete la retta $z+u=h+d$, ma con velocità non costante. La realizzazione è possibile per $s<h$: quando l'ombra si avvicina alla base, il riflettore deve allontanarsi indefinitamente. Non è possibile proseguire oltre $s=h$ con un riflettore finito sulla parete e l'obelisco fermo. La traiettoria riflessa rispetto a $y=0$ è analoga.

\paragraph{Punto (c).}
Fissiamo $L=(-d,0,h+d)$ e spostiamo la base in $P=(x,y,0)$, con punta $(x,y,h)$. L'intersezione della semiretta luminosa con il terreno è
\[
Q=\left(h+\frac{h+d}{d}x,\frac{h+d}{d}y,0\right).
\]
Imponendo $Q=(h-s,s,0)$ si ottiene la traiettoria indicata nella risposta. La base si muove quindi lungo una retta, a velocità costante; in particolare riproduce tutto il tratto realizzabile nel punto precedente senza richiedere velocità divergenti.
}
